% /root/monodromy with suggestions by /root, 19 November 2026.
\subsection{Finite constant data for an arbitrary perfect base field}
\label{subsec:perfect-field-monodromy-data}
\status{Proved; independent parent proof and original-source audit passed.
The output is finite polynomial data and matrix lists, rather than a
uniform structural classification of the resulting finite groups.}

\begin{lemma}\label{lem:effective-separable-factorization}
Let $k_0$ be a finitely generated field over $\mathbf F_p$,
given by a finite field presentation.  Arithmetic and
factorization of separable univariate polynomials over
$k_0(t_1,\ldots,t_M)$ are terminating exact procedures.
For a polynomial that defines a finite separable field
extension of this rational function field, its absolutely
irreducible factors over $\bar k_0(t_1,\ldots,t_M)$ can
also be found by a terminating procedure, with their
coefficients in explicitly presented finite separable
extensions of $k_0$.
\end{lemma}

\begin{proof}
We give details so that no absolute-factorization or
Galois-group oracle is concealed in the assertion.

A separating transcendence basis of $k_0/\mathbf F_p$
and a finite separable primitive presentation over the
corresponding rational function field can be found
by elimination, linear algebra and enumeration.
Explicitly, enumerate candidate rational functions
for a transcendence basis, check algebraic independence
by elimination, and check that the remaining algebra
is finite by a zero-dimensional Groebner computation.
Separability is checked by the rank of the differentials,
or by the nonzero derivative of a primitive minimal
polynomial.  Such a candidate exists because the
ground field is perfect.  For a finite separable
extension, the primitive element is found by testing
linear combinations of a finite basis, as in the
finite-grid argument of the previous theorem.
Thus we may write
\[
 k_0(t)=R(\alpha),\qquad
 R=\mathbf F_q(u_1,\ldots,u_s,t_1,\ldots,t_M),
\]
with a separable minimal polynomial for $\alpha$.
All arithmetic is then polynomial reduction and linear
algebra over a rational function field.

Here is a direct factorization procedure over $R(\alpha)$.
For a monic separable polynomial $f(Z)$, try translates
$g_c(Z)=f(Z+c\alpha)$ for successively enumerated
$c\in R$, until
\[
 N_c(Z)=\operatorname{Norm}_{R(\alpha)/R}(g_c(Z))
\]
is square-free.  Its norm is computed by a resultant.
Only finitely many $c$ are bad: over an algebraic
closure its roots are $\beta-c\sigma(\alpha)$, where
$\beta$ ranges over the roots of $f^\sigma$ and
$\sigma$ over the distinct embeddings of $R(\alpha)$.
Within one embedding they are distinct; across two
different embeddings an equality excludes at most
one value of $c$.  The rational function field $R$
is infinite here.  If there are no transcendental
variables, the field is finite and bounded polynomial
enumeration directly factors $f$ instead.

Factor $N_c$ over $R$ by clearing denominators and
the bounded finite-field polynomial enumeration in
Step 4 of Theorem~\ref{thm:effective-monodromy-matrices}.
For each irreducible factor $h$ of $N_c$, compute
$\gcd(g_c,h)$ over $R(\alpha)$ and discard the units.
These gcds are irreducible and give the factors of
$g_c$.  To verify irreducibility, two roots of $g_c$
in the same $R$-Galois orbit are carried to one another
by an automorphism fixing $R(\alpha)$: otherwise the
target root would occur in two different conjugates
of $g_c$, contradicting square-freeness of $N_c$.
Undoing the translate factors $f$.  Thus all separable
factorization steps used below are effective.

For absolute factorization, clear denominators of the
given separable polynomial and remove its content
in $k_0[t]$, obtaining a primitive polynomial
$F(t,Z)$.  Its total degree bounds the total degree
of every polynomial factor over $\bar k_0$.
For each possible pair of leading monomials and
factor degrees, write two polynomials with unknown
coefficients, normalize the leading coefficient of
one to one, and equate their product with $F$.
Impose the selected $Z$-degree of the proposed factor
to be positive and require the selected leading
coefficients to be nonzero (an inverse variable
imposes this as an equation).  Permit the complementary
factor to be the unit polynomial; in particular the
full-degree factor is always one of the cases.
There are finitely many cases.  The coefficient
ideal is proper if and only if that case has a
factorization over the algebraic closure, by the
Nullstellensatz.  Properness is decided by whether
its Groebner basis contains one.

Select the least positive $Z$-degree for which a
factor exists.  A factor of that degree is absolutely
irreducible: a nontrivial splitting would yield a
smaller positive $Z$-degree.  There are no factors
independent of $Z$, by primitivity and Gauss's lemma.
Thus a finite list of coefficient tests proves what
degree to seek; it is not a search whose negative
outcome is undecidable.

We can find the factor's coefficients explicitly.
The field defined by the original irreducible factor
is separably generated over $k_0$.  Its algebraic
constant field is finite separable, and after adjoining
it its components are geometrically integral.
Indeed finiteness and geometric irreducibility are
\cite[Lemmas~10.47.8 and 10.47.13]{Stacks037P,Stacks037Q};
geometric reducedness follows from separability by
\cite[Lemmas~10.44.2 and 10.44.3]{Stacks030W,Stacks030X}.
Consequently the normalized absolute factors have
coefficients in finite separable extensions of $k_0$.
Enumerate separable irreducible polynomials over
$k_0$, testing irreducibility by the procedure above,
and enumerate coefficient tuples in their finite
extensions.  Dovetail these enumerations and test
the finitely many coefficient equations.  A solution
exists in one such extension, so the search terminates.
Repeat division and the finite coefficient tests
until all factors have been found; their positive
$Z$-degrees sum to the original degree.
This proves the absolute-factorization assertion
without assuming a preexisting algebraic closure
computing oracle.
\end{proof}

\begin{theorem}\label{thm:perfect-field-exact-monodromy-data}
Let $k$ be any perfect field of characteristic $p$,
and let the fixed smooth projective embedding $X/k$
be defined over a finitely generated subfield
$k_0\subset k$.  Fix $d\geq d_{\rm coh}$.
There is a terminating construction from a finite
presentation of this field of definition and the
equations of $X$ of the following finite data:
\[
 H_{\rm geom}\triangleleft H_0\subset\GL_{r_d}(\mathbf F_p),
 \qquad E/k_0\text{ finite Galois},\qquad
 \rho:H_0\twoheadrightarrow\operatorname{Gal}(E/k_0),
 \quad\ker\rho=H_{\rm geom}.
\]
The two groups are output as complete lists of matrices,
$E$ by a finite field presentation, and $\rho$ by its
action on generators of $E$.
In a compatible geometric frame, the requested images
over $k$ are exactly
\begin{equation}\label{eq:perfect-field-monodromy-formula}
 G_{\rm geom}=H_{\rm geom},\qquad
 G_{\rm arith}=\rho^{-1}
       \bigl(\operatorname{Gal}(E/(E\cap k))\bigr).
\end{equation}
The intersection is taken using the chosen embedding
of $E$ into a separable algebraic closure of $k_0$
inside an algebraic closure of $k$.
Changing the initial frame conjugates all the displayed
matrix groups simultaneously.

Geometric irreducibility is a finite subspace test on
$H_{\rm geom}$.  Arithmetic irreducibility is the same
test on the explicit preimage in
\eqref{eq:perfect-field-monodromy-formula}; the only
additional base-field information is the finite
intersection $E\cap k$.
\end{theorem}

\begin{proof}
The cohomological matrix and stable-image algorithm
of Proposition~\ref{prop:effective-generic-rank} work
over $k_0(t)$ as well: they use rational linear
algebra and Frobenius linearization, not perfection
of this intermediate field.  Form the frame algebra
and find a primitive element as before.
Use Lemma~\ref{lem:effective-separable-factorization}
to factor its polynomial over $k_0(t)$ and choose
one field component $L$.
The finite list of matrix substitutions preserving
that component is $H_0$.
It acts as the full Galois group of $L/k_0(t)$:
the component is a connected torsor under its
stabilizer, whose order equals its degree.

Let $E$ be the algebraic closure of $k_0$ in $L$.
It is finite separable by the field facts cited
in the preceding lemma, and it is Galois over $k_0$
because $L/k_0(t)$ is Galois.
The field $L/E$ is separably generated with $E$
algebraically closed in it, hence geometrically
integral by the same cited lemmas.
It follows that the components after extending
constants to $\bar k_0$ are indexed by the embeddings
of $E$; their stabilizer in $H_0$ is
\[
 H_{\rm geom}=\operatorname{Gal}(L/E(t)).
\]
The restriction map to $E$ is $\rho$.
Both matrix lists are obtained by the finite ideal
preservation tests, now using one absolute factor
for $H_{\rm geom}$.

We explain how the field $E$ itself is recovered,
rather than leaving it as an unspecified relative
algebraic closure.  Let $f_0$ be the primitive
polynomial of $L$, and choose a monic absolutely
irreducible factor $f_{00}$ over $\bar k_0(t)$.
Normalize every rational-function coefficient by
coprime numerator and denominator with one leading
coefficient equal to one.
Let $E_0$ be the field over $k_0$ generated by all
the scalar coefficients in these normalized expressions.
The factorization algorithm gives them in a finite
explicit extension, so $E_0$ is computed by adjoining
them and closing their vector span under multiplication.
This terminates inside that finite extension.

In fact $E_0=E$ with the selected embedding.
Since $L/E$ is geometrically integral, the geometric
components are indexed by the embeddings of $E$.
The stabilizer of one component under
$\operatorname{Gal}(\bar k_0/k_0)$ is exactly
$\operatorname{Gal}(\bar k_0/E)$.
On the other hand, that automorphism fixes the
chosen factor precisely when it fixes all of its
normalized coefficients, so this stabilizer is
$\operatorname{Gal}(\bar k_0/E_0)$.
The finite separable fields are equal.
Thus $E$ has been explicitly obtained.

Its embedding into $L$ and the restriction map
$\rho$ can also be calculated by linear algebra.
The selected factor realizes $L$ inside
$E(t)[Z]/(f_{00})$.
This latter field has degree
$[E:k_0]\deg f_{00}=[L:k_0(t)]$ over $k_0(t)$,
so it is the same field $L$.
Express the computed generators of $E$ in the
basis $1,\theta,\ldots,\theta^{[L:k_0(t)]-1}$
by solving a linear system.  Apply every retained
matrix substitution to these expressions and
reduce.  This outputs $\rho$ explicitly.

Finally extend the constants to $k$.
Put $F=E\cap k$.  Since $E/k_0$ is finite Galois,
$\operatorname{Gal}(Ek/k)=\operatorname{Gal}(E/F)$.
Since $L/E$ is geometrically integral, extending
$E$ to $Ek$ does not change its geometric degree
or create new components.  The component over
$k(t)$ therefore has degree
\[
 |H_{\rm geom}|[E:F],
\]
and its automorphism group is exactly the preimage
of $\operatorname{Gal}(E/F)$ under $\rho$.
This proves \eqref{eq:perfect-field-monodromy-formula}.
The geometric image is unchanged by any extension
of the algebraically closed constant field, giving
the first formula.  Subspace enumeration gives the
irreducibility assertions exactly as before.
\end{proof}

\begin{remark}\label{rem:perfect-field-data-scope}
The theorem is an exact finite-data description for
an arbitrary perfect field, and an algorithm whenever
the finite field-of-definition presentation and the
finite constant intersection can be evaluated.
It makes no claim that membership in an arbitrary
unpresented, possibly noncomputable perfect field
can be decided by a universal program.
For a finite base field, the intersection is computed
by finite-field arithmetic and the previous theorem
already gives a complete unconditional algorithm.
The structural question of identifying these outputs
uniformly as familiar finite groups is separate;
for projective space it was answered above, whereas
for general ambient varieties it is not inferred
from this construction.
\end{remark}
